% A.87 reserved by the root.
\subsection{Dyadic inertia fixes every odd multiplier}

\begin{statement}{A.87}{Proved}
{A dyadic symplectic constraint and recognition of thirty-six}
For $t\in\Q_2\setminus\{0,1\}$, define its contact order $h(t)>0$
by the unique one of
\[
 (v_2(t),v_2(1-t))=(h(t),0),\quad(0,h(t)),
                         \quad(-h(t),-h(t)).
\tag{A.87.1}
\]
Every Legendre curve with $h(t)=2$ has minimal discriminant
valuation eight and inertial group $\mathrm{SL}_2(\F_3)$ of
order twenty-four.  Every Legendre curve with $h(t)=4$
has inertial group of order at most two.

Consider a labeled integral realization $E_{\lambda}$ in A.14
with $h(\lambda)=2$ at two.  Its ordinary dyadic parameter
$x_2$ has $h(x_2)=2$ as well.  For the fixed global
Weil multipliers of A.82 one has
\[
                    \Delta_\ell\in(\Z_\ell^\times)^2
                              \qquad\text{for every odd }\ell.
\tag{A.87.2}
\]
This conclusion makes no assumption on $\Delta_2$.
For an even integer $\lambda=4u$, $u$ odd, the actual
dyadic valuation pair is exactly $(2,0)$ whenever
\[
 u\equiv1\text{ or }3\pmod8,
             \qquad\text{or}\qquad
 \lambda\equiv0\text{ or }1\pmod9.
\tag{A.87.2a}
\]
Consequently, within their original parameter ranges, the criteria of A.79 and
A.82 apply to these parameters with an arbitrary labeled
integral realization, without their additional full-$H$
restriction at dyadic valuation two.
In particular, every parameter $\lambda=3^n+1$, $n\ge1$,
is recognized with an arbitrary labeled integral realization.

Every noncuspidal Selmer section with labeled integral
realization $E_{36}$ is the point section $s_{36}$.
\end{statement}
\begin{proof}
We give the local inertia calculation before applying it to
the global multipliers.  Write
\[
 E_t:\ y^2=x(x-1)(x-t),\qquad
 E_t^{(d)}:\ y^2=x(x-d)(x-dt).
\tag{A.87.3}
\]
The latter is the quadratic twist specified by the actual
isomorphism $(x,y)\mapsto(dx,d\sqrt d\,y)$ over
$\Q_2(\sqrt d)$.  Two coordinate identities used below are
\[
 E_t\xrightarrow{\ (x,y)\mapsto(x-1,y)\ }
                E_{1-t}^{(-1)},
       \qquad E_t=E_{1/t}^{(t)}.
\tag{A.87.4}
\]
The second equality is equality of the displayed cubic
equations, with the last two roots interchanged.  These maps
specify the twists at the other two labels.

First suppose $t=4u$, $u\in\Z_2^\times$.
The displayed integral model has
\[
\begin{aligned}
 c_4&=16(16u^2-4u+1),\\
 c_6&=64(4u+1)(8u-1)(2u-1),\\
 \Delta&=2^8u^2(1-4u)^2.
\end{aligned}
\tag{A.87.5}
\]
Thus its invariant valuations are $(4,6,8)$, and it is
minimal because a nonminimal integral equation would lower
the discriminant valuation by a positive multiple of twelve.
Its normalized discriminant satisfies
$\widetilde\Delta\equiv1\pmod8$.
At contact one, write $t=1-4u$; the same formulas have the
same $c_4,\Delta$ and the opposite $c_6$, so the same
valuation and discriminant assertions hold.
At contact infinity write $t=u/4$.  The actual change
$X=4x$, $Y=8y$ gives the integral equation
\[
 Y^2=X(X-4)(X-u),\qquad
 \begin{aligned}
 c_4&=16(u^2-4u+16),\\
 c_6&=64(u+4)(u-2)(u-8),\\
 \Delta&=2^8u^2(u-4)^2.
 \end{aligned}
\tag{A.87.6}
\]
Again the valuations are $(4,6,8)$ and
$\widetilde\Delta\equiv1\pmod8$.

Here is a faithful finite inertia computation, including
the coefficient-three reduction.  Since $v_2(j)=4$, each
curve in (A.87.5)--(A.87.6) has potentially good reduction.
The finite inertia image on $T_3(E)$ is faithful on $E[3]$:
the reduction kernel $1+3M_2(\Z_3)$ is torsion-free.
Indeed, for $M=1+3^aN$, $a\ge1$, with at least one unit
entry of $N$, the binomial formula gives
$v_3(M^m-1)=a+v_3(m)$, where matrix valuation is the
minimum entry valuation.  No positive power can therefore
be the identity.
The Weil pairing makes the inertia determinant one because
$\mu_3$ is unramified over $\Q_2$.
The inertial group $\Phi$ consequently embeds in
$\mathrm{SL}_2(\F_3)$, which has order twenty-four.
The field cut out by the kernel of this inertia representation
over $\Q_2^{\mathrm{un}}$ is the minimal field of good reduction;
equivalently it is
$\Q_2^{\mathrm{un}}(E[3])$, as in
\cite[\S4.1]{LegFreitasKraus2022}.

Put $e=|\Phi|$.  Good reduction over this ramification
degree changes discriminant valuation $8e$ to zero by a
twelfth power.  Hence $12\mid8e$ and $3\mid e$.
There is no subgroup of order twelve in
$\mathrm{SL}_2(\F_3)$: an index-two subgroup would be normal
and give a nontrivial homomorphism to $\Z/2$, whereas
\[
 \begin{pmatrix}1&1\\0&1\end{pmatrix},\qquad
 \begin{pmatrix}1&0\\1&1\end{pmatrix}
\tag{A.87.7}
\]
generate the group and both have order three.
Thus $e\in\{3,6,24\}$.
The necessary minimal-invariant conditions for $e=3$ in
\cite[Table~5, case $C_{3,\mathrm{iii}}$]{LegFreitasKraus2022}
and the complete $e=6$ list in
\cite[\S14.3]{LegFreitasKraus2022} both require
$\widetilde\Delta\equiv-1\pmod4$ for the valuation triple
$(4,6,8)$.  Our normalized discriminants are one modulo
four.  Both orders are excluded, proving $e=24$.
The embedding is therefore onto
$\mathrm{SL}_2(\F_3)$.

For comparison, let $t=16u$, $u\in\Z_2^\times$.
On the twist $E_{16u}^{(-1)}$, the change
$x=4X$, $y=8Y+4X$ gives exactly
\[
               Y^2+XY=X^3+4uX^2+uX.
\tag{A.87.8}
\]
Its discriminant is $u^2(1-16u)^2$, a unit, so this
curve has good reduction over $\Q_2$.
Thus $E_{16u}$ becomes good over a quadratic extension.
For contact one, (A.87.4) identifies $E_{1-16u}$ with
this same good twist.  For contact infinity,
$E_t=E_{1/t}^{(t)}$ and $1/t=16/u$; this curve is the
quadratic twist by $-t$ of the good curve
$E_{16/u}^{(-1)}$.
These exact maps prove $e\le2$ at all three contact-four
labels.

Return to the given global realization with
$h(\lambda)=2$.  Its local section at two is ordinary
by A.37, since the realized curve is potentially good.
The same representation makes its actual point curve
$E_{x_2}$ potentially good.
The two labeled Kummer squareclasses in A.16 force both
entries of its valuation pair to be even.
Potential good reduction allows only contact orders
one through four; thus only two or four remain.
The actual integral intertwiner identifies the finite
inertia images.  The orders twenty-four and at most two
computed above exclude contact four.
Consequently both $E_{\lambda}$ and $E_{x_2}$ have
minimal discriminant valuation eight and the same
inertial field.
For any odd coefficient $\ell$, this equality of fields
also follows directly by taking the kernels of the
isomorphic inertia actions on $T_\ell$; finite inertia
reduction modulo $\ell$ is faithful by the same
torsion-free-kernel argument.

Freitas--Kraus's criterion
\cite[Theorem~7, with the numbering of arXiv v5]{LegFreitasKraus2022}
says that for common inertial field of order twenty-four,
the isomorphisms on odd torsion have square Weil
multiplier when $(2/\ell)=1$, and, when
$(2/\ell)=-1$, exactly when the two minimal discriminant
valuations agree modulo three.  It also says that the
two multiplier squareclasses cannot both occur.
Here the two valuations are equal to eight, so the
reduction of every actual $T_\ell$ intertwiner has square
multiplier.  Square units lift from $\F_\ell$ to
$\Z_\ell$ by Hensel's lemma.
In passing from the restricted global section to its
ordinary local point, its geometric representative
change has Legendre matrix in $H$ and determinant one.
The Weil-paired fiber identifications likewise preserve
the pairing.  Thus the multiplier just computed is
the same fixed $\Delta_\ell$ of the global comparison,
proving (A.87.2).  This argument makes no assertion
about the coefficient two.

For $\lambda=4u$ with $u\equiv1$ or $3\pmod8$,
the prescribed dyadic squareclasses of $(A,B)$ are
$(u,5)$, after removing the square $4$ from $A$.
A contact-one point of order two has $x_2\equiv5\pmod8$,
which contradicts the first class.
A contact-infinity point $x_2=v/4$ with $v\in\Z_2^\times$
would have the unit squareclasses $(v,4-v)$.
For $v\equiv1$ or $3\pmod8$ its second class is three
or one, respectively, and cannot be five.
Therefore the dyadic pair is $(2,0)$ in these two unit classes.

We prove the second condition in (A.87.2a) and keep all
weights in the coefficient-three congruence.
For any $\lambda=4u$, the possible dyadic pairs are now
$(d_A,d_B)=(2,0),(0,2),(-2,-2)$.
Write $f_p=v_p(\lambda)$ for odd $p\mid\lambda$ and
$e_q=v_q(1-\lambda)$ for odd $q\mid1-\lambda$.
The exact global support expansions are
\[
 \begin{aligned}
 A&=\operatorname{sgn}(\lambda)\,
                    2^{d_A}\prod_{p\mid\lambda,\ p\text{ odd}}p^{\alpha_p},\\
 B&=\operatorname{sgn}(1-\lambda)\,
                    2^{d_B}\prod_{q\mid1-\lambda}q^{\beta_q}.
 \end{aligned}
\tag{A.87.2b}
\]
All odd bad primes retain their realized contact labels.
The integral Tate comparison and (A.87.2) at coefficient
three give
\[
 \alpha_{p,3}/f_p\in(\Z_3^\times)^2,
       \qquad \beta_{q,3}/e_q\in(\Z_3^\times)^2.
\tag{A.87.2c}
\]
Every square unit is one modulo three.  Hence
$\alpha_{p,3}-f_p\in3f_p\Z_3$ and
$\beta_{q,3}-e_q\in3e_q\Z_3$.
Since $\log_3(r)\in3\Z_3$ for every rational prime
$r\ne3$ and $\log_3(3)=0$, these assertions imply
\[
 \begin{aligned}
 \sum_p\alpha_{p,3}\log_3p
                      &\equiv\log_3(\lambda/4)\pmod9,\\
 \sum_q\beta_{q,3}\log_3q
                      &\equiv\log_3(1-\lambda)\pmod9.
 \end{aligned}
\tag{A.87.2d}
\]
The signs have zero logarithm.  In particular these
congruences remain valid when some weights $f_p$ or
$e_q$ are divisible by three; their divisibility then
only increases the error valuation.

Suppose first that $\lambda\equiv1\pmod9$ and the
dyadic pair is one of the two alternatives.
A cusp at three must have label one.  Its regular
coordinate $A_3=1$, together with (A.87.2d), would give
\[
       0\equiv(d_A-2)\log_3(2)+\log_3(\lambda)
                  \equiv(d_A-2)\log_3(2)\pmod9.
\tag{A.87.2e}
\]
Here $d_A=0$ or $-2$, whereas
$\log_3(2)\equiv6\pmod9$ has valuation one.
Neither $-2$ nor $-4$ is divisible by three, so this
is impossible.
If instead $\lambda\equiv0\pmod9$, a cusp at three
would have label zero and regular coordinate $B_3=1$.
The same calculation gives
\[
       0\equiv d_B\log_3(2)+\log_3(1-\lambda)
                      \equiv d_B\log_3(2)\pmod9,
\tag{A.87.2f}
\]
again impossible because $d_B=2$ or $-2$.

Thus, under either congruence and either alternate
dyadic pair, the restriction at the potentially
multiplicative prime three is ordinary.
A.34 applied to the restricted global intertwiners
makes their entire multiplier square, including its
coefficient-two component.  The explicit scalar
adjustment in A.17 therefore gives an $H$-conjugator.
The two separate modulo-eight characters
$(A_8,B_8)$ of $H^{\mathrm{ab}}$ in A.56 now give
\[
                      (d_A,d_B)\equiv(2,0)\pmod8.
\tag{A.87.2g}
\]
Both $(0,2)$ and $(-2,-2)$ violate this congruence.
This contradiction excludes the alternate pairs and
proves the second condition in (A.87.2a).

The pair $(2,0)$ is precisely the support conclusion for which
A.79 and A.82 used their full-$H$ assumption at $k=2$;
their remaining proofs require only the labeled
integral comparison.  The stated extension of those
criteria follows.
For $\lambda=3^{2j+1}+1$, the dyadic valuation is two.
The case $j=0$ is $\lambda=4$, with $u=1$, and lies
in the first condition of (A.87.2a).
For $j\ge1$ one has $\lambda\equiv1\pmod9$.
In every case the odd neighbour has the single prime
factor three, so every minimum in A.79 is unique;
that criterion proves recognition without an $H$
assumption.  The even exponents $n$ in $3^n+1$ have
dyadic valuation one and are already covered by A.41.

We finish with $\lambda=36=4\cdot9$.
At the odd potentially multiplicative primes $3,5,7$,
the labels are respectively zero, one, one, with
contact orders two, one, one.
The dyadic pair is $(2,0)$ by the preceding calculation.
Signs, the finite-support Kummer map and the full
integral Tate valuation comparison give
\[
 A=4\cdot3^{\alpha},\qquad B=-5^{\beta}7^{\gamma},
 \qquad X=\alpha/2,\ \beta,\gamma\in\widehat\Z^\times.
\tag{A.87.9}
\]
At ordinary bad primes the same statement follows from
the exact local parameter in A.78.
Keeping the same global multiplier as above, the
paired local Tate comparisons give
\[
 X_\ell=\frac{a_{3,\ell}^2}{\Delta_\ell},\qquad
 \beta_\ell=\frac{a_{5,\ell}^2}{\Delta_\ell},\qquad
 \gamma_\ell=\frac{a_{7,\ell}^2}{\Delta_\ell},
       \qquad a_{q,\ell}\in\Z_\ell^\times.
\tag{A.87.10}
\]
In particular all three quantities are square units
at every odd coefficient by (A.87.2).

A cusp at five would have label one and regular
coordinate $A_5=1$.  Its coefficient-five logarithm
would give
\[
          X_5=-\frac{\log_5(2)}{\log_5(3)}
                      \equiv-2/4=2\pmod5.
\tag{A.87.11}
\]
Here $\log_5(2)\equiv10$ and
$\log_5(3)\equiv20\pmod{25}$, so both logarithms
have exact valuation one and division is valid.
The residue two is a nonsquare, contradicting
(A.87.10).  Thus the local section at five is
ordinary.  A.78 fixes its exact value at $36$ and
therefore fixes the full global exponent $\beta=1$.

A cusp at seven would also have label one.
Its actual ordinary Tate parameter and formal cusp
Tate class are, respectively,
\[
 q_7=\frac{(-35)^2}{256}F(-35),\qquad
          \frac{B^2}{256}
                  =\frac{5^2 7^{2\gamma}}{256}.
\tag{A.87.12}
\]
These representations are split at seven: the nodal
slopes at contact one are $\pm1$ on both sides.
Their valuation ratio is $\gamma_7$.
The full coefficient-seven Tate extension equality
of A.63, with $\log_7(7)=0$, is consequently
\[
 2\log_7(5)-8\log_7(2)
     =\gamma_7\bigl(2\log_7(5)-8\log_7(2)
                                  +\log_7 F(-35)\bigr).
\tag{A.87.13}
\]
The integral inverse-nome series and its logarithm
give $\log_7 F(-35)\equiv-35\pmod{49}$ by
(A.84.5).  The exact normalized logarithms are
\[
       \log_7(5)/7\equiv1,
                   \qquad \log_7(2)/7\equiv5\pmod7.
\tag{A.87.14}
\]
Thus, after division by seven, the numerator in
(A.87.13) is $2-40=4\pmod7$ and its denominator
is $4-5=6\pmod7$.  This denominator is a unit,
and the equality forces
\[
                         \gamma_7\equiv4/6=3\pmod7.
\tag{A.87.15}
\]
Three is a nonsquare modulo seven, again contradicting
(A.87.10).  Hence no label-one cusp remains.
The only possible label is zero at three, so A.37
proves that the full section is a rational point.
The labeled primitive-isogeny argument of A.54 fixes
that point at $36$.
\end{proof}

The finite group, invariant, twist and logarithm checks
are retained in
\path{code/legendre_dyadic_symplectic_certificate.py};
their complete output is saved beside the program.
The general finite-support recognition problem remains
open; (A.87.2) supplies a proved restriction on its
odd-coefficient multipliers.
